% ============================================================
% MAIN.TEX — SuperMarket Chain DataBase
% Stile basato su template tesi UniUD (Edoardo Scudeller)
% Compilare con: pdflatex -> bibtex -> pdflatex -> pdflatex
% ============================================================

\documentclass[12pt, a4paper, twoside, openright]{article}

% ------------------------------------------------------------
% ENCODING E LINGUA
% ------------------------------------------------------------
\usepackage[T1]{fontenc}
\usepackage[utf8]{inputenc}
\usepackage[italian]{babel}

% ------------------------------------------------------------
% LAYOUT E MARGINI
% ------------------------------------------------------------
\usepackage[top=2.5cm, bottom=2.5cm, left=3cm, right=2.5cm]{geometry}
\usepackage{setspace}
\onehalfspacing

% ------------------------------------------------------------
% FONT
% ------------------------------------------------------------
\usepackage{lmodern}

% ------------------------------------------------------------
% GRAFICA E COLORI
% ------------------------------------------------------------
\usepackage{graphicx}
\usepackage[dvipsnames, table]{xcolor}
\graphicspath{{immagini/}}

% ------------------------------------------------------------
% TABELLE
% ------------------------------------------------------------
\usepackage{booktabs}
\usepackage{tabularx}
\usepackage{multirow}
\usepackage{longtable}
\usepackage{array}

% ------------------------------------------------------------
% MATEMATICA
% ------------------------------------------------------------
\usepackage{amsmath}
\usepackage{amssymb}
\usepackage{amsthm}

% ------------------------------------------------------------
% PSEUDOCODICE E ALGORITMI
% ------------------------------------------------------------
\usepackage{algorithm}
\usepackage{algpseudocode}          % stile algorithmicx (algpseudocode)
\algrenewcommand\algorithmicindent{1.5em}
\algnewcommand\algorithmicand{\textbf{ and }}
\algnewcommand\algorithmicor{\textbf{ or }}
\algnewcommand\algorithmicnot{\textbf{not }}
\algnewcommand\algorithmictrue{\textbf{true}}
\algnewcommand\algorithmicfalse{\textbf{false}}
% Rinomina "Require" e "Ensure" in italiano:
\renewcommand{\algorithmicrequire}{\textbf{Input:}}
\renewcommand{\algorithmicensure}{\textbf{Output:}}
% Rinomina la parola "Algorithm" nel float:
\floatname{algorithm}{Algoritmo}

% ------------------------------------------------------------
% LISTATI DI CODICE SQL
% ------------------------------------------------------------
\usepackage{listings}
\lstdefinestyle{sqlstyle}{
    language=SQL,
    basicstyle=\ttfamily\small,
    keywordstyle=\color{Blue}\bfseries,
    commentstyle=\color{OliveGreen}\itshape,
    stringstyle=\color{BrickRed},
    numberstyle=\tiny\color{gray},
    numbers=left,
    stepnumber=1,
    numbersep=8pt,
    frame=single,
    frameround=tttt,
    backgroundcolor=\color{gray!6},
    breaklines=true,
    showstringspaces=false,
    tabsize=4,
    captionpos=b
}
\lstset{style=sqlstyle}

% ------------------------------------------------------------
% HYPERLINK E PDF METADATA
% ------------------------------------------------------------
\usepackage[hidelinks, pdfusetitle]{
hyperref}
\hypersetup{
    pdftitle  = {SuperMarket Chain DataBase},
    pdfauthor = {-- INSERIRE NOMI --},
    pdfsubject= {Progetto Basi di Dati},
    pdfkeywords = {database, SQL, PostgreSQL, supermarket, relational}
}

% ------------------------------------------------------------
% INTESTAZIONI E PIÈ DI PAGINA
% ------------------------------------------------------------
\usepackage{fancyhdr}
\pagestyle{fancy}
\fancyhf{}
\fancyhead[LE,RO]{\thepage}
\fancyhead[LO]{\nouppercase{\rightmark}}
\fancyhead[RE]{\nouppercase{\leftmark}}
\renewcommand{\headrulewidth}{0.4pt}

% ------------------------------------------------------------
% INDICE CON MINITOC PER SEZIONE
% ------------------------------------------------------------
\usepackage{minitoc}
\setcounter{tocdepth}{3}
\setcounter{secnumdepth}{3}

% ------------------------------------------------------------
% RIFERIMENTI E BIBLIOGRAFIA
% ------------------------------------------------------------
\usepackage[numbers, sort&compress]{natbib}
\bibliographystyle{plainnat}

% ------------------------------------------------------------
% ALTRI PACCHETTI UTILI
% ------------------------------------------------------------
\usepackage{multicol}       % layout multicolonna (usato nel frontespizio)
\usepackage{enumitem}       % personalizzazione liste
\usepackage{caption}        % personalizzazione didascalie
\usepackage{subcaption}     % sotto-figure
\usepackage{float}          % controllo posizionamento float [H]
\usepackage{tikz}           % disegni/diagrammi
\usetikzlibrary{er, positioning, arrows.meta, shapes}
\usepackage{pdflscape}      % pagine orizzontali con \begin{landscape}

% ============================================================
% INIZIO DOCUMENTO
% ============================================================
\begin{document}

\dominitoc  % inizializza minitoc

% ------------------------------------------------------------
% FRONTESPIZIO
% ------------------------------------------------------------
% ============================================================
% FRONTESPIZIO - PERSONALIZZARE
% Inserire qui: nome università, corso, titolo, relatore, laureando
% ============================================================

\thispagestyle{empty}
\begin{center}

% LOGO UNIVERSITÀ (sostituire con il proprio logo nella cartella frontespizio/)
\includegraphics[width=4cm]{frontespizio/logo.pdf} \\

% UNIVERSITA - NOME
{\LARGE \textbf{Università degli Studi di -- INSERIRE --}} \\

% LINEA DIVISORIA
\vspace{0.15cm}
\rule{0.2\textwidth}{0.4pt}\\
\vspace{0.65cm}

% FACOLTA - NOME
{\Large {-- Dipartimento / Facoltà --}} \\
\vspace{1cm}

% CORSO DI LAUREA - NOME
{\Large {-- Corso di Laurea --}} \\
\vspace{3cm}

% TIPO DI ELABORATO E TITOLO
{\Large {Progetto di Basi di Dati}} \\
\vspace{0.25cm}
{\Large \textbf{SuperMarket Chain: Progettazione e}}\\ [0.1cm]
{\Large \textbf{Implementazione di un Database Relazionale}}
\end{center}
\vspace{4cm}

% DOCENTE & STUDENTI
\begin{multicols}{2}
\raggedright
{\large Docente:}\\[0.1cm]
{\large \textbf{Prof. -- NOME COGNOME --}}

\vspace{0.5cm}

% Eventuale secondo docente:
% {\large {Correlatore:}}\\[0.1cm]
% {\large \textbf{-- NOME --}}

\columnbreak

\raggedleft
{\large Studente/i:}\\[0.1cm]
{\large \textbf{-- Nome Cognome --}}\\
{\large \textbf{-- Nome Cognome --}}\\
% {\large {Matricola: ------}}
\end{multicols}

\begin{center}

% LINEA DIVISORIA
\vspace{1.65cm}
\rule{0.8\textwidth}{0.4pt}\\
\vspace{0.15cm}

% ANNO ACCADEMICO
\large{\textbf{Anno Accademico -- / --}}

\end{center}
\newpage


% ------------------------------------------------------------
% SOMMARIO
% ------------------------------------------------------------
\begin{abstract}
\noindent
Il presente progetto affronta la progettazione e l'implementazione di un database relazionale
per la gestione di una catena di supermercati (\textit{SuperMarket Chain}).
L'elaborato descrive l'analisi dei requisiti, la progettazione concettuale tramite schema E/R,
la traduzione al modello relazionale, la normalizzazione e l'implementazione fisica in PostgreSQL.
Vengono inoltre presentate le principali query SQL per l'interrogazione e la gestione dei dati,
corredati di pseudocodice per gli algoritmi applicativi rilevanti.
\end{abstract}

\newpage

% ------------------------------------------------------------
% INDICE GENERALE
% ------------------------------------------------------------
\tableofcontents
\newpage

\listoffigures
\listoftables
\newpage

% ============================================================
% SEZIONE 0 — INTRODUZIONE
% ============================================================
\section*{Introduzione}
\addcontentsline{toc}{section}{Introduzione}

Il progetto ha come obiettivo la realizzazione di un sistema di gestione dati per una catena
di supermercati distribuita su più sedi. Il sistema deve supportare la gestione di prodotti,
fornitori, dipendenti, clienti, ordini e inventario, garantendo integrità referenziale
e performance adeguate ai volumi operativi tipici della grande distribuzione organizzata (GDO).

\newpage

% ============================================================
% SEZIONE 1 — ANALISI DEI REQUISITI
% ============================================================
\section{Analisi dei Requisiti}
\minitoc
\vspace{1cm}

\subsection{Descrizione del dominio applicativo}

% TODO: Descrivere il dominio (catena supermercati, sedi, prodotti, ecc.)

\subsection{Requisiti funzionali}

% TODO: Elencare i requisiti funzionali del sistema

\subsection{Requisiti non funzionali}

% TODO: Performance, disponibilità, sicurezza, scalabilità

\newpage

% ============================================================
% SEZIONE 2 — PROGETTAZIONE CONCETTUALE (SCHEMA E/R)
% ============================================================
\section{Progettazione Concettuale}
\minitoc
\vspace{1cm}

\subsection{Schema Entità-Relazione}

% Inserire qui il diagramma E/R come immagine:
% \begin{figure}[H]
%     \centering
%     \includegraphics[width=\textwidth]{immagini/schema-er.pdf}
%     \caption{Schema Entità-Relazione del database SuperMarket Chain.}
%     \label{fig:schema-er}
% \end{figure}

\subsection{Dizionario dei dati}

\subsubsection{Entità principali}

% TODO: tabella con entità, attributi, chiavi

\subsubsection{Relazioni principali}

% TODO: tabella con relazioni, cardinalità, partecipazioni

\newpage

% ============================================================
% SEZIONE 3 — PROGETTAZIONE LOGICA (MODELLO RELAZIONALE)
% ============================================================
\section{Progettazione Logica}
\minitoc
\vspace{1cm}

\subsection{Traduzione E/R $\rightarrow$ Relazionale}

% TODO: Schemi relazionali con chiavi primarie sottolineate e FK indicate

\subsection{Normalizzazione}

% TODO: verifica 1NF, 2NF, 3NF / BCNF

\newpage

% ============================================================
% SEZIONE 4 — IMPLEMENTAZIONE FISICA (PostgreSQL)
% ============================================================
\section{Implementazione Fisica in PostgreSQL}
\minitoc
\vspace{1cm}

\subsection{Creazione delle tabelle}

% Esempio di listato SQL:
% \begin{lstlisting}[caption={Creazione tabella PRODOTTO.}, label={lst:create-prodotto}]
% CREATE TABLE Prodotto (
%     codice_prodotto  SERIAL       PRIMARY KEY,
%     nome             VARCHAR(100) NOT NULL,
%     categoria        VARCHAR(50),
%     prezzo_unitario  NUMERIC(8,2) NOT NULL CHECK (prezzo_unitario >= 0),
%     codice_fornitore INTEGER      REFERENCES Fornitore(codice_fornitore)
% );
% \end{lstlisting}

\subsection{Vincoli e trigger}

% TODO: trigger di integrità, vincoli CHECK, DEFAULT

\subsection{Indici e ottimizzazione}

% TODO: indici B-Tree, ottimizzazione query

\newpage

% ============================================================
% SEZIONE 5 — QUERY SQL PRINCIPALI
% ============================================================
\section{Query SQL Principali}
\minitoc
\vspace{1cm}

\subsection{Interrogazioni di base}

% TODO: SELECT, JOIN, WHERE di base

\subsection{Interrogazioni avanzate}

% TODO: aggregazioni, subquery, window functions, CTE

\newpage

% ============================================================
% SEZIONE 6 — ALGORITMI E PSEUDOCODICE
% ============================================================
\section{Algoritmi Applicativi e Pseudocodice}
\minitoc
\vspace{1cm}

% ------------------------------------------------------------
% ESEMPIO 1: Algoritmo di rifornimento automatico
% ------------------------------------------------------------
\subsection{Rifornimento Automatico dell'Inventario}

L'algoritmo di rifornimento monitora le scorte di ogni prodotto per sede e genera
automaticamente un ordine di rifornimento quando la giacenza scende sotto la soglia minima.

\begin{algorithm}[H]
\caption{Rifornimento Automatico Inventario}
\label{alg:rifornimento}
\begin{algorithmic}[1]
\Require Lista delle sedi $S$, soglia minima $\theta$, quantità di riordino $Q_r$
\Ensure Ordini di rifornimento generati per tutti i prodotti sotto soglia

\ForAll{$s \in S$} \Comment{Per ogni sede della catena}
    \ForAll{prodotto $p$ in $s$}
        \State $q \gets$ \Call{GetGiacenza}{$s$, $p$}
        \If{$q < \theta_p$}
            \State $\text{ord} \gets$ \Call{CreaOrdine}{$s$, $p$, $Q_r$}
            \State \Call{InviaAlFornitore}{$\text{ord}$, $p.\text{fornitore}$}
            \State \Call{LogOrdine}{$\text{ord}$}
        \EndIf
    \EndFor
\EndFor
\end{algorithmic}
\end{algorithm}

% ------------------------------------------------------------
% ESEMPIO 2: Ricerca del fornitore ottimale
% ------------------------------------------------------------
\subsection{Selezione del Fornitore Ottimale}

Data una richiesta di approvvigionamento per un prodotto, l'algoritmo seleziona il fornitore
che minimizza il costo totale (prezzo unitario + costo di trasporto).

\begin{algorithm}[H]
\caption{Selezione Fornitore a Costo Minimo}
\label{alg:fornitore}
\begin{algorithmic}[1]
\Require Prodotto $p$, quantità richiesta $q$, insieme fornitori $F_p$
\Ensure Fornitore ottimale $f^*$

\State $c_{\min} \gets +\infty$
\State $f^* \gets \textbf{null}$
\ForAll{$f \in F_p$}
    \State $c_f \gets f.\text{prezzoUnitario} \times q + f.\text{costoTrasporto}$
    \If{$c_f < c_{\min}$}
        \State $c_{\min} \gets c_f$
        \State $f^* \gets f$
    \EndIf
\EndFor
\State \Return $f^*$
\end{algorithmic}
\end{algorithm}

% ------------------------------------------------------------
% ESEMPIO 3: Calcolo del totale di uno scontrino (cassa)
% ------------------------------------------------------------
\subsection{Calcolo Totale Scontrino}

\begin{algorithm}[H]
\caption{Calcolo Totale Scontrino con Sconti}
\label{alg:scontrino}
\begin{algorithmic}[1]
\Require Lista articoli $A = \{(p_i, q_i)\}$, cliente $c$
\Ensure Totale da pagare $T$

\State $T \gets 0$
\ForAll{$(p_i, q_i) \in A$}
    \State $\text{prezzo} \gets p_i.\text{prezzoUnitario} \times q_i$
    \If{\Call{IsInPromozione}{$p_i$}}
        \State $\text{prezzo} \gets \text{prezzo} \times (1 - p_i.\text{sconto})$
    \EndIf
    \State $T \gets T + \text{prezzo}$
\EndFor
\If{\Call{IsFidelityCard}{$c$}}
    \State $T \gets T \times (1 - \text{SCONTO\_FIDELITY})$
\EndIf
\State \Return $\Call{Round}{T, 2}$
\end{algorithmic}
\end{algorithm}

\newpage

% ============================================================
% SEZIONE 7 — CONCLUSIONI
% ============================================================
\section*{Conclusioni}
\addcontentsline{toc}{section}{Conclusioni}

% TODO: Sintesi dei risultati, limitazioni, sviluppi futuri

\newpage

% ============================================================
% APPENDICE
% ============================================================
\appendix
\section{Script SQL Completo}

% TODO: script DDL completo

\section{Dati di Test}

% TODO: script di popolamento (INSERT)

\newpage

% ============================================================
% BIBLIOGRAFIA
% ============================================================
\bibliography{bibliografia}

\end{document}
